\chapter{Этаж~2. Составной узел: кварки и квантовые числа}
\label{ch:floor2}

\section{Радиус и объект}
\label{sec:floor2-scale}

Этаж~2 на оси длины (\bcref{prop:ladder-scale-steps}):
\[
  N\sim 10^{15}\cdot\caStepSpace,\qquad \log_{10}N\approx 15.
\]
Это уже \emph{составной узел}: несколько фокусов планконов в одной $\varepsilon$-окрестности,
связанных топологией \latGCKshort\ и запретом Паули (A16), а не один вихрь на одном~$\hv$.

На~$\layerMacro$ здесь --- схема конфайнмента и кварковых степеней свободы;
на~$\layerMicro$ --- полный учёт устойчивых составных конфигураций с нетривиальным
топологическим и цветовым бюджетом.

\section{Конфайнмент без нового поля}
\label{sec:floor2-confinement}

«Вытащить кварк'' на макроязыке --- разорвать составной топологический узел.
На~$\Mlayer$ это не извлечение одного планкона из ячейки, а
\emph{пара} (или больше) планконов с разнесёнными фокусами:
энергия растёт с расстоянием, потому что нарушается голономия и занятость на звезде~$\varepsilon$.

Цветовой заряд $SU(3)$ \emph{не вводится} отдельной калибровкой сверху.
Кратность устойчивых фокусов в барионном или мезонном узле
выводится непосредственно из переписи полных подграфов (клик) на \latGCKshort,
запрета Паули и модуля~$K_P$.

\begin{theorem}[Перепись клик $K_m$ и три ортогональные плоскости цвета на \latGCKshort]
\label{thm:floor2-fcc-cliques}
Пусть $K_m$ --- полный подграф из $m$ узлов~$\Lambda$, в котором любая пара вершин
соединена ребром ближайших соседей длины $1\cdot\caStepSpace$ (прямой каузальный обмен за один такт~$\htick$ по A1).
Тогда вокруг каждого узла $x_0\in\Lambda$ на \latGCKshort\ ($\geomNeighborTwelve=12$):
\begin{enumerate}[noitemsep]
\item \textbf{$m=2$ ($K_2$, двухфокусный димер):} ровно $12$ рёбер (мезонный канал $B=0$ / пара $\pm$);
\item \textbf{$m=3$ ($K_3$, минимальный замкнутый цикл, $\mathrm{girth}=3=d$):} ровно $24$ равносторонних треугольника
  (барионный узел $B=1$ из $d=3$ взаимно связанных фокусов);
\item \textbf{$m=4$ ($K_4$, максимальная клика):} ровно $8$ правильных тетраэдров над $8$ треугольными гранями кубооктаэдра;
  при этом каждый треугольник~$K_3$ достраивается до тетраэдра~$K_4$ ровно \emph{одной} четвёртой вершиной в~$\Lambda$
  (зеркальная вершина исключена условием чётности ГЦК $x_1+x_2+x_3\equiv 0\pmod 2$);
\item \textbf{$m\ge 5$ ($K_{\ge 5}$):} строго $0$ клик --- компактный однотактовый узел из пяти и более взаимных соседей
  на \latGCKshort\ геометрически невозможен.
\end{enumerate}
У любого треугольника~$K_3$ (например, $\{(0,0,0),(0,1,1),(1,0,1)\}$) три его ребра
\[
  e_{12}=(0,1,1)\in YZ,\qquad
  e_{13}=(1,0,1)\in XZ,\qquad
  e_{23}=(1,-1,0)\in XY
\]
лежат в \emph{трёх взаимно ортогональных координатных плоскостях} $\{YZ,XZ,XY\}$.
\end{theorem}

\begin{proof}
Из определения~\latGCKshort\ $12$ соседей начального узла имеют вид перестановок $(\pm 1,\pm 1,0)$ с квадратом длины~$2$.
Два соседа $u,v\in N_{12}$ образуют с началом координат правильный треугольник~$K_3$ тогда и только тогда, когда $|u-v|^2=2$:
каждая из $12$ вершин кубооктаэдра имеет степень~$4$ на его рёбрах, что даёт $12\times 4/2=24$ треугольника~$K_3$.
Разности координат трёх вершин любого такого треугольника обнуляют по очереди каждую из трёх осей ($x_1=0$, $x_2=0$, $x_3=0$),
поэтому три ребра~$K_3$ принадлежат трём взаимно перпендикулярным плоскостям $YZ$, $XZ$ и $XY$.
Общий сосед всех трёх вершин $K_3=\{(0,0,0),(0,1,1),(1,0,1)\}$ в~$N_{12}$ единственен --- узел $(1,1,0)$,
что даёт $24/3=8$ тетраэдров~$K_4$ и ноль полных подграфов при $m\ge 5$.
\end{proof}

Следствие~\cref{thm:floor2-fcc-cliques} снимает мистику с «трёх цветов'' кварков:
в одной координатной плоскости два вихря противоположного знака $(+1,-1)$ складывают фазовые градиенты в противофазе
и аннигилируют (этаж~1), тогда как на треугольнике~$K_3$ три фокуса закручены в трёх взаимно перпендикулярных плоскостях
($XY\perp YZ\perp XZ$). «Цвет'' фокуса --- это его координатная плоскость вращения на ребре~$K_3$,
а «бесцветный синглет'' бариона --- замкнутый треугольник~$K_3$, охватывающий все три ортогональные плоскости $d=3$.

\section{Квантовые числа, происхождение «третей'' заряда и магнитные моменты нуклонов}
\label{sec:floor2-quantum-numbers}

На этаже~2 впервые появляются составные квантовые числа,
которые на этаже~0 заданы лишь внутренними степенями одного планкона:

\begin{table}[htbp]
  \centering
  \caption{Откуда берутся квантовые числа составного узла на этаже~2.}
  \label{tab:floor2-qnumbers}
  \footnotesize
  \setlength{\tabcolsep}{4pt}
  \begin{tabular}{@{}lll@{}}
    \toprule
    число & этаж~0 (один $\hv$) & этаж~2 (клика $K_3$ / $K_2$ на \latGCKshort) \\
    \midrule
    электрический $Q$ & целый $n\in\mathbb{Z}$ (A10) & целый $Q\in\mathbb{Z}$ на узле; на~$T$ делится на $d=3$ вершины ($\pm 1/3,\pm 2/3$) \\
    спин & $SU(2)$ на ядре, $2\pi\to -1$ & $\tfrac12\otimes\tfrac12\otimes\tfrac12 = \tfrac12\ (p,n)\oplus\tfrac32\ (\Delta)$ \\
    цвет ($N_c=3$) & --- & $3$ ортогональные плоскости $\{XY,YZ,XZ\}$ рёбер $K_3$ (\cref{thm:floor2-fcc-cliques}) \\
    поколение ($N_{\mathrm{gen}}=3$) & --- & $d=3$ пространственных стека упаковки (гл.~\ref{ch:si-sm}) \\
    барионное $B$ & $b\in\{0,1\}$ на ячейке & замкнутый треугольный замок $K_3$ ($\mathrm{girth}=3$, $\Delta B=0$) \\
    \bottomrule
  \end{tabular}
\end{table}

\paragraph{Почему на~$\Mlayer$ нет дробных зарядов, а на~$\layerMacro$ видны $+2/3$ и $-1/3$.}
По аксиоме~A10 топологический заряд каждой ячейки строго целый ($n\in\mathbb{Z}$).
Однако в барионном треугольнике~$K_3$ единичный целый заряд $Q=+1$ (протон) или $Q=0$ (нейтрон)
распределён по $d=3$ геометрически равноправным вершинам, которые по запрету Паули (A16)
разбиты по спину в отношении $2:1$ ($\uparrow\uparrow\downarrow$).
Макроскопический прибор на~$\layerMacro$, делящий целый заряд треугольника на $d=3$ вершины в пропорции $2:1$
с единичным шагом перезарядки $\Delta Q=\pm 1$, приписывает вершинам эффективные партонные веса
\[
  q_u^{(\mathrm{eff})}=+\frac{2}{d}=+\frac{2}{3},\qquad
  q_d^{(\mathrm{eff})}=-\frac{1}{d}=-\frac{1}{3},\qquad
  2q_u^{(\mathrm{eff})}+q_d^{(\mathrm{eff})}=+1,\quad
  q_u^{(\mathrm{eff})}+2q_d^{(\mathrm{eff})}=0.
\]
Тем самым «дробные заряды кварков'' --- не нарушение целочисленности~A10 внутри ячейки~$\hv$,
а проекция целого заряда на три вершины треугольной плакетки~$K_3$.

\paragraph{Магнитные моменты протона и нейтрона и «загадка спина протона''.}
В формуле массы протона (\bcref{ch:si-sm}) уже зафиксирован геометрический квант упаковки однотактового кубооктаэдра
на одно ребро звезды:
\[
  \delta_{\mathrm{pack}} = \frac{\caKappaGeom^2}{\geomNeighborTwelve} = \frac{1/2}{12} = \frac{1}{24}.
\]
Каждая из $d=3$ вершин барионного треугольника~$K_3$ отдаёт в упаковку оболочки один квант~$\delta_{\mathrm{pack}}$,
поэтому суммарная доля потерь на трёх вершинах равна обратному числу каналов иерархии $\vevHierarchySteps=8$:
\[
  d\cdot\delta_{\mathrm{pack}} = 3\times\frac{1}{24} = \frac{1}{8} = \frac{1}{\vevHierarchySteps}.
\]
Отсюда без единого подгоночного параметра следуют изоскалярный момент нуклона ($\mu_p+\mu_n$, объясняющий,
почему спиновый вклад фокусов равен $1-1/8=7/8$, а не~$1$) и индивидуальные магнитные моменты протона
(вычет $d_{\mathrm{out}}(E)=5$ рёберных квантов~$\delta_{\mathrm{pack}}$) и нейтрона
(компенсация $k=2$ квантов~$\delta_{\mathrm{pack}}$ из $m_n=m_p+2m_e$):
\begin{align}
  \mu_p + \mu_n &= \bigl(1 - d\,\delta_{\mathrm{pack}}\bigr)\mu_N
    = \frac{7}{8}\,\mu_N = 0{,}8750\,\mu_N
    \quad(\text{PDG: }0{,}8798\,\mu_N,\ \text{откл.\ }0{,}55\%),\\
  \mu_p &= \bigl(d - 5\,\delta_{\mathrm{pack}}\bigr)\mu_N
    = \frac{67}{24}\,\mu_N \approx +2{,}7917\,\mu_N
    \quad(\text{PDG: }+2{,}7928\,\mu_N,\ \text{откл.\ }0{,}042\%),\\
  \mu_n &= \bigl(-2 + 2\,\delta_{\mathrm{pack}}\bigr)\mu_N
    = -\frac{23}{12}\,\mu_N \approx -1{,}9167\,\mu_N
    \quad(\text{PDG: }-1{,}9130\,\mu_N,\ \text{откл.\ }0{,}189\%).
\end{align}

\section{Составной 13-ячеечный электрон и лестница кварковых масс в нативных единицах \texorpdfstring{$(dh,dt)$}{(dh,dt)} и \texorpdfstring{$(s_0,\ladderEnergy,m_{\arg})$}{(s0,E0,marg)}}
\label{sec:floor2-quark-masses-natural}

На носителе~$\Mlayer$ внешние координаты --- счётчики $(dh,dt)=(\caStepSpace,\htick)$ с тактовой скоростью
$\caSignalSpeed=dh/dt$, а вся механика порождается одной ступенью действия~$\ladderAction$ на $(\hv,1\,dt)$
(\bcref{ch:descent,sec:coords-dh-dt}):
\[
  \ladderMomentum=\frac{\ladderAction}{dh},\qquad
  \ladderEnergy=\frac{\ladderAction}{dt},\qquad
  \ladderForce=\frac{\ladderAction}{dh\cdot dt},\qquad
  m_{\arg}=\frac{2\ladderMomentum}{\caSignalSpeed}=\frac{2\ladderAction\,dt}{dh^2}.
\]
Учебниковые $E_P=\ladderEnergy/\caKappaGeom$ и $m_P=m_{\arg}/\caKappaGeom$
живут только в мосте $\Mlayer\to\layerMacro$, где $\caKappaGeom=R_{\mathrm{in}}/R_{\mathrm{out}}$ ---
отношение вписанной и описанной сфер кубооктаэдра, а $r=\Delta\varphi_{\min}/(2\pi)$
(на~$\Mlayer$: $\Delta\varphi_{\mathrm{disc}}/N_{\mathrm{ring}}$) --- фазовый остаток планковской дырки ($\Omega=1/r$).
Поскольку гауссов множитель моста равен $\caKappaGeom/\sqrt{r}$, при переходе от $E_P=\ladderEnergy/\caKappaGeom$
к нативным единицам ячейки $(\ladderEnergy,m_{\arg})$ фактор~$\caKappaGeom$ в~$v$ сокращается тождественно,
и все формулы записываются без голых корней только через $(\alpha_{\mathrm{fs}},\caKappaGeom,r,\vevHierarchySteps,N_\varphi,\geomNeighborTwelve)$:
\[
  v = \frac{1}{\sqrt{r}}\,\alpha_{\mathrm{fs}}^{\vevHierarchySteps}\,\ladderEnergy,\qquad
  m_H^{(\mathrm{bare})} = \caKappaGeom^2\,v = \frac{\caKappaGeom^2}{\sqrt{r}}\,\alpha_{\mathrm{fs}}^{\vevHierarchySteps}\,m_{\arg}.
\]

\begin{enumerate}[leftmargin=*]
\item \textbf{Электрон --- составной $13$-ячеечный кластер ($1+12$).}
  В силу теоремы о минимальном ореоле $R_{\mathrm{dress}}=1\cdot dh$ (этаж~1, \cref{prop:floor1-dress-radius})
  единичный заряд $Q=\pm 1$ не может быть точечным: вокруг центрального ядра ($1$ узел) принудительно возбуждена
  вся $\varepsilon$-звезда первых соседей ($\geomNeighborTwelve=12$ узлов).
  Минимальный одетый электрон образует $13$-ячеечный кластер кубооктаэдра $N_{\mathrm{cluster}}=1+\geomNeighborTwelve=13=N_\varphi$,
  поэтому его масса в единицах~$m_{\arg}$ делится на~$N_\varphi$:
  \[
    \massElectron = \frac{\alpha_{\mathrm{fs}}^2\,m_H}{N_\varphi},\qquad
    \massElectron^{(\mathrm{bare})} = \frac{\caKappaGeom^2}{N_\varphi\sqrt{r}}\,\alpha_{\mathrm{fs}}^{\vevHierarchySteps+2}\,m_{\arg}.
  \]

\item \textbf{Тяжёлые фокусы бариона против квадратичных зарядовых шагов ($u,d$).}
  Свободные кварки не наблюдаются; в таблицах сосуществуют два масштаба:
  \begin{itemize}[noitemsep]
  \item \emph{структурная (конституэнтная) масса фокуса на треугольнике~$K_3$:}
    каждая из $d=3$ вершин протона несёт треть полной массы узла,
    \[
      M_q^{(\mathrm{focus})} = \frac{m_p}{d} = \frac{\caKappaGeom^2}{d\sqrt{r}}\,\alpha_{\mathrm{fs}}^{\vevHierarchySteps+1}\left(1+\frac{\caKappaGeom^2}{\geomNeighborTwelve}\right)m_{\arg};
    \]
  \item \emph{токовые ($\overline{\mathrm{MS}}$) массы $m_u,m_d$:}
    электромагнитная цена перезарядки вершин треугольника~$K_3$ в единицах минимального заряженного кванта канала~$\massElectron$
    (как в $m_n-m_p=2\massElectron$). При числителе зарядового веса $u$-вершины $n_u=2$ ($+2/3$) и полном числе вершин
    треугольника $d=3$ для $d$-вершины квадратичные вклады равны целым квадратам:
    \[
      m_u = n_u^2\,\massElectron = 4\,\massElectron,\qquad
      m_d = d^2\,\massElectron = 9\,\massElectron,\qquad
      \frac{m_u}{m_d} = \frac{n_u^2}{d^2} = \frac{4}{9}.
    \]
  \end{itemize}

\item \textbf{Старшие кварковые ступени ($s,c,b,t$) через $(\caKappaGeom,r,\alpha_{\mathrm{fs}})$ в единицах $m_{\arg}$.}
  Четыре тяжёлые кварковые ступени задаются без подгоняемых констант Юкавы и без голых корней через~$m_{\arg}$, $m_p$,
  коэффициент квадратичного роста сфер кубооктаэдра $\geomNeighborTwelve-2=10$,
  отношение числа вершин тетраэдра~$K_4$ к треугольнику~$K_3$ ($(d+1)/d$),
  мягкую грань $(1-\caKappaGeom^6)/\caKappaGeom^6=N_4+n_{\mathrm{SU}(2)}$ и геометрический фактор кубооктаэдра $\caKappaGeom=R_{\mathrm{in}}/R_{\mathrm{out}}$:
  \begin{align}
    m_s &= \frac{m_p}{\geomNeighborTwelve-2},\\
    m_c &= \frac{d+1}{d}\,m_p,\\
    m_b &= \frac{1-\caKappaGeom^6}{d\,\caKappaGeom^6}\,\alpha_{\mathrm{fs}}\,v
      = \frac{1-\caKappaGeom^6}{d\,\caKappaGeom^6\sqrt{r}}\,\alpha_{\mathrm{fs}}^{\vevHierarchySteps+1}\,m_{\arg},\\
    m_t^{(\mathrm{bare})} &= \caKappaGeom\,v = \frac{m_H^{(\mathrm{bare})}}{\caKappaGeom}
      = \frac{\caKappaGeom}{\sqrt{r}}\,\alpha_{\mathrm{fs}}^{\vevHierarchySteps}\,m_{\arg}.
  \end{align}
  Тем самым загадка Стандартной модели «почему константа Юкавы топ-кварка равна единице''
  получает прямой геометрический ответ: на полке~$v$ бозон Хиггса отвечает квадрату отношения сфер кубооктаэдра
  $\caKappaGeom^2$ ($m_H^{(\mathrm{bare})}=\caKappaGeom^2 v$),
  а верхний потолок одиночного фермионного фокуса (топ-кварк) --- первой степени $\caKappaGeom$
  ($m_t^{(\mathrm{bare})}=\caKappaGeom v$).

\item \textbf{Три поколения заряженных лептонов ($e,\mu,\tau$) из $(d=3,\caKappaGeom,\lfloor\caBitBudget\rfloor)$.}
  Три пространственных стека поколений ($g\in\{0,1,2\}$, $N_{\mathrm{gen}}=d=3$) образуют симметричную тройку
  с фазовым шагом $2\pi g/d$ вокруг главной диагонали, амплитудой $\caKappaGeom^{-1}=R_{\mathrm{out}}/R_{\mathrm{in}}$ кубооктаэдра
  и начальным фазовым сдвигом двухкомпонентного спинора~$\mathbb{C}^2$ на битовом регистре ячейки
  $\delta_0 = 2/\lfloor\caBitBudget\rfloor = 2/9$:
  \[
    \sqrt{m_g} \propto 1 + \frac{1}{\caKappaGeom}\cos\!\left(\frac{2\pi g}{d} + \frac{2}{\lfloor\caBitBudget\rfloor}\right),
    \qquad g\in\{0,1,2\}.
  \]
  Поскольку $\caKappaGeom^2=1/2$, отношение суммы масс к квадрату суммы амплитуд (тождество Коидэ)
  выполняется тождественно:
  \[
    \frac{\massElectron + m_\mu + m_\tau}{\bigl(\sqrt{\massElectron}+\sqrt{m_\mu}+\sqrt{m_\tau}\bigr)^2}
    = \frac{1 + (2\caKappaGeom^2)^{-1}}{d} = \frac{d-1}{d} = \frac{2}{3},
  \]
  а отношения масс поколений без единого подгоняемого параметра равны
  \begin{align*}
    \frac{m_\mu}{\massElectron} &\approx 206{,}7703\quad(\text{PDG: }206{,}7683,\ \text{откл.\ }+9{,}83\ \mathrm{ppm}),\\
    \frac{m_\tau}{\massElectron} &\approx 3477{,}47\quad(\text{PDG: }3477{,}23\pm 0{,}23,\ \text{внутри }1{,}05\sigma).
  \end{align*}
\end{enumerate}

\section{Связь с сильным взаимодействием, атом водорода (\texorpdfstring{$k=9$}{k=9}) и упаковка \texorpdfstring{$N_{\mathrm{pack}}\in\mathbb{Z}$}{Npack in Z} (\texorpdfstring{$k=10$}{k=10})}
\label{sec:floor2-strong}

На этаже~2 сильное взаимодействие --- не отдельный Лагранжиан,
а \emph{короткий} канал упаковки и удержания той же фазовой башни,
что даёт электрослабое сопряжение на больших~$N$
(\bcref{ch:alpha}, \bcref{ch:si-sm}):
$\alpha_s(v)=d/(\vevHierarchySteps\pi)=3/(8\pi)$.
На 3D кипящем океане \texttt{VACUUM\_BOIL} растяжение одной вершины трёхплоскостного треугольника~$K_3$
на расстояние $r_{\mathrm{pull}}=1,2,3,4\cdot dh$ монотонно увеличивает число надпороговых узлов облака~$\rho_\Theta$
($+1017\to +1048\to +1195\to +1311$ ячеек) и суммарное фазовое натяжение $\sum|\Delta\varphi|$
($4565\to 4666\to 4830\to 5130$), демонстрируя геометрический конфайнмент без введения калибровочного поля~$SU(3)$.

Шкала семени иерархии~$\alpha$ смыкает геометрические сферы этажа~1 ($N_{12}^m$, $k\le 5$)
с инфракрасной hop-лестницей ($k\ge 7$), атомом водорода ($k=9$) и упаковкой второго кристалла ($k=10$):
\begin{enumerate}[noitemsep]
\item один шаг иерархии $\alpha_{\mathrm{fs}}^{-1}\approx 137{,}036\approx \geomNeighborTwelve^2=144$ отвечает двум координационным сферам ГЦК;
\item восемь битов иерархии ($\vevHierarchySteps=\lfloor\caBitBudget\rfloor-1=8$) дают электрослабый масштаб
  $v=\frac{1}{\sqrt{r}}\alpha_{\mathrm{fs}}^{\vevHierarchySteps}\ladderEnergy$ ($\lambda_v/dh \approx 4{,}95\times 10^{16}$, $k=6$);
\item девятая ступень ($\alpha_{\mathrm{fs}}^{\vevHierarchySteps+1}$, барионная упаковка $B=1$ на~$K_3$) фиксирует комптоновский шаг протона
  $N_p = m_P/m_p \approx 1{,}304\times 10^{19}$;
\item десятая ступень ($\alpha_{\mathrm{fs}}^{\vevHierarchySteps+2}$ с фазовым бюджетом $N_\varphi=13$) выходит прямо на комптоновский якорь электрона
  $N_c = m_P/\massElectron \approx 2{,}379\times 10^{22}$ ($k=8$) и классический радиус $N_{-1}=\alpha_{\mathrm{fs}}N_c$ ($k=7$);
\item \textbf{одиннадцатая ступень ($\alpha_{\mathrm{fs}}^{\vevHierarchySteps+d}=\alpha_{\mathrm{fs}}^{11}$, $k=9$) --- атом водорода~$H$ и энергия связи в единицах $\ladderEnergy,m_{\arg}$:}
  целое число шагов первой боровской оболочки вокруг барионного ядра $B=1$ равно
  $N_{a0}=\mathrm{round}(N_c/\alpha_{\mathrm{fs}})\in\mathbb{Z}$ ($\approx 3{,}259\times 10^{24}$ шагов~$dh$),
  а энергия связи (дефект массы $\Delta m_H$) по решёточному закону Кулона и теореме о вириале равна
  \[
    \frac{\Delta m_H}{m_{\arg}} = \frac{BE_H}{\ladderEnergy} = \frac{\alpha_{\mathrm{fs}}}{2\caKappaGeom\,N_{a0}} \approx \frac{1}{2\hat M\,N_{a0}}
  \]
  ($BE_H\approx 13{,}667\ \mathrm{эВ}$ против $13{,}606\ \mathrm{эВ}$ CODATA, откл.\ $0{,}45\%$),
  то есть связь протона и электрона в атоме водорода --- один фазовый квант $\ladderEnergy=\ladderAction/dt$,
  распределённый по $\hat M=97$ местам силы и $2N_{a0}$ шагам диаметра орбиты;
\item \textbf{двенадцатая ступень ($k=10$) --- точная целочисленность упаковки этажей $N_{\mathrm{pack}}\in\mathbb{Z}$:}
  поскольку и планковская ячейка Вороного $\mathrm{d}V=\caKappaGeom\,dh^3$, и атомная ячейка Вигнера--Зейтца второго кристалла
  с целым ребром $A=K\cdot dh$ ($K\in\mathbb{N}$) имеют одну и ту же форму ромбического додекаэдра ГЦК,
  иррациональный фактор~$\caKappaGeom$ сокращается тождественно:
  \[
    N_{\mathrm{pack}}(K) = \frac{V_{\mathrm{WS}}(K\cdot dh)}{\mathrm{d}V}
    = \frac{\caKappaGeom\,(K\cdot dh)^3}{\caKappaGeom\,dh^3}
    = K^3 \in \mathbb{Z}
    \qquad\bigl(\text{при }K=N_{a0}:\ N_{\mathrm{pack}}=N_{a0}^3\in\mathbb{Z}\bigr).
  \]
\end{enumerate}

\section{Статус этажа~2}
\label{sec:floor2-open}

\paragraph{Закрыто (алгебра, симуляция растяжения $K_3$ и сверка с PDG).}
\begin{itemize}[noitemsep]
\item перепись клик $K_2=12$, $K_3=24$, $K_4=8$, $K_{\ge 5}=0$ и единственность кирального апекса $K_3\to K_4$ (\cref{thm:floor2-fcc-cliques});
\item три ортогональные плоскости $\{XY,YZ,XZ\}$ рёбер $K_3$ как геометрическая природа трёх цветов ($N_c=d=3$) и монотонный рост объёма облака~$\rho_\Theta$ при растяжении $r_{\mathrm{pull}}=1\ldots 4$;
\item происхождение эффективных весов $+2/3$ и $-1/3$ на~$\layerMacro$ из деления целого заряда треугольника~$K_3$ на $d=3$ вершины;
\item магнитные моменты $\mu_p=67/24\,\mu_N$, $\mu_n=-23/12\,\mu_N$ и $\mu_p+\mu_n=7/8\,\mu_N$ через квант упаковки $\delta_{\mathrm{pack}}=\caKappaGeom^2/\geomNeighborTwelve=1/24$;
\item лестница кварковых ступеней ($u,d,s,c,b,t$) и трёх поколений лептонов ($e,\mu,\tau$) при $\delta_0=2/\lfloor\caBitBudget\rfloor$ ($+9{,}83\ \mathrm{ppm}$, $+69{,}8\ \mathrm{ppm}$);
\item полные параметры матрицы смешивания кварков CKM ($\lambda = 35/156 \approx 0{,}22436$, $A = (\geomNeighborTwelve-2)/\geomNeighborTwelve = 5/6 \approx 0{,}8333$, гауссова факторизация $\hat M = 4^2+9^2 = 97 \implies \bar\rho + i\bar\eta = (4+9i)/(2N_\varphi) = (4+9i)/26$, $\gamma = \arctan(9/4) \approx 66{,}04^\circ$) и матрицы PMNS ($\sin^2\theta_{12}=4/13$, $\sin^2\theta_{23}=7/13$, $\sin^2\theta_{13}=1/46$) выведены без подгоночных углов в \bcref{ch:si-sm} (\bcref{prop:ckm-pmns-sectors});
\item непрерывная стыковка шкалы $k=5\to 10$, энергия связи атома водорода $BE_H/\ladderEnergy=\alpha_{\mathrm{fs}}/(2\caKappaGeom N_{a0})$ и целочисленная упаковка $N_{\mathrm{pack}}(K)=K^3\in\mathbb{Z}$.
\end{itemize}

После этажа~2 лестница выходит на инфракрасную ось
$N_k=N_c/\alpha^{k}$ (\bcref{eq:ir-hop-ladder}): там начинается атомная и макро-физика на~$\layerMacro$.

\paragraph{Что читать дальше.}
Константы связи и массы --- \bcref{ch:alpha}, \bcref{ch:si-sm};
гравитация и электрослабый сектор как макроскопическое считывание той же среды --- в главах о силах модели
(см. оглавление).
